\documentclass[letterpaper,11pt]{article}

\usepackage{latexsym}
\usepackage[empty]{fullpage}
\usepackage{titlesec}
\usepackage{marvosym}
\usepackage[usenames,dvipsnames]{color}
\usepackage{enumitem}
\usepackage[hidelinks]{hyperref}
\usepackage{fancyhdr}
\usepackage[portuguese]{babel}
\usepackage{tabularx}
\usepackage{fontawesome5}
\usepackage{multicol}
\setlength{\multicolsep}{-3.0pt}
\setlength{\columnsep}{-1pt}
\input{glyphtounicode}
\hypersetup{
  pdftitle={Kaylon Souza - Currículo},
  pdfauthor={Kaylon Souza},
  pdfsubject={Currículo técnico},
  pdfkeywords={Kaylon Souza, currículo, personal page, Node.js, Python, TypeScript, React Native, Flask, JavaScript, software}
}

\pagestyle{fancy}
\fancyhf{}
\fancyfoot{}
\renewcommand{\headrulewidth}{0pt}
\renewcommand{\footrulewidth}{0pt}

\addtolength{\oddsidemargin}{-0.6in}
\addtolength{\evensidemargin}{-0.5in}
\addtolength{\textwidth}{1.19in}
\addtolength{\topmargin}{-.7in}
\addtolength{\textheight}{1.4in}

\urlstyle{same}
\raggedbottom
\raggedright
\setlength{\tabcolsep}{0in}

\titleformat{\section}{
  \vspace{-4pt}\scshape\raggedright\large\bfseries
}{}{0em}{}[\color{black}\titlerule \vspace{-5pt}]

\pdfgentounicode=1

\newcommand{\resumeItem}[1]{
  \item\small{{#1 \vspace{-2pt}}}
}
\newcommand{\resumeSubheading}[4]{
  \vspace{-2pt}\item
    \begin{tabular*}{1.0\textwidth}[t]{l@{\extracolsep{\fill}}r}
      \textbf{#1} & \textbf{\small #2} \\
      \textit{\small #3} & \textit{\small #4} \\
    \end{tabular*}\vspace{-7pt}
}
\newcommand{\resumeProjectHeading}[2]{
    \item
    \begin{tabular*}{1.001\textwidth}{l@{\extracolsep{\fill}}r}
      \small#1 & \textbf{\small #2}\\
    \end{tabular*}\vspace{-7pt}
}
\newcommand{\resumeSubHeadingListStart}{\begin{itemize}[leftmargin=0.0in, label={}]}
\newcommand{\resumeSubHeadingListEnd}{\end{itemize}}
\newcommand{\resumeItemListStart}{\begin{itemize}}
\newcommand{\resumeItemListEnd}{\end{itemize}\vspace{-5pt}}

\begin{document}

\begin{center}
    {\Huge \scshape Kaylon Souza} \\ \vspace{1pt}
    Macajuba - BA, Brasil \\ \vspace{1pt}
    \small \raisebox{-0.1\height}\faPhone\ (11) 91659-1346 ~
        \href{mailto:kaylon.contact@gmail.com}{\raisebox{-0.2\height}\faEnvelope\ \underline{kaylon.contact@gmail.com}} \\
    \href{https://kaaylooon.github.io}{\raisebox{-0.2\height}\faGlobe\ \underline{kaaylooon.github.io}} ~
        \href{https://linkedin.com/in/kaylonsouza}{\raisebox{-0.2\height}\faLinkedin\ \underline{linkedin.com/in/kaylonsouza}} ~
        \href{https://github.com/kaaylooon}{\raisebox{-0.2\height}\faGithub\ \underline{github.com/kaaylooon}}
    \vspace{-8pt}
\end{center}


\section{Resumo}
Full-stack | Node.js • TypeScript • React Native • Python, com experiência prática em desenvolvimento de aplicações próprias.

\section{Formação}
\resumeSubHeadingListStart
  \resumeSubheading
      {Colégio Estadual Carlito de Carvalho}{Fev. 2024 -- Dez. 2026}
      {Curso Técnico de Administração Integrado ao Ensino Médio}{Bahia, Brasil}
\resumeSubHeadingListEnd

\section{Projetos}
\resumeSubHeadingListStart
\resumeProjectHeading{\href{https://kaaylooon.github.io/projetos/motriz/}{\textbf{Motriz}} | \emph{TypeScript, Expo, React Native, AsyncStorage}}{2026 -- Presente}
\resumeItemListStart
  \resumeItem{Aplicativo mobile para organizar ciclos de estudo e tarefas, combinando distribuição ponderada de disciplinas com uma lista inspirada em todo.txt.}
  \resumeItem{Gera blocos de estudo a partir de pesos por dificuldade e quantidade de tópicos, com intercalação para reduzir repetições próximas.}
  \resumeItem{Mantém persistência local de ciclo, disciplinas, tarefas e configurações usando AsyncStorage e tipagem estrita com TypeScript.}
\resumeItemListEnd

\resumeProjectHeading{\href{https://kaaylooon.github.io/projetos/katalogo/}{\textbf{Katálogo}} | \emph{Python, Flask, SQLite/PostgreSQL, HTML/CSS}}{2024}
\resumeItemListStart
  \resumeItem{Aplicação web completa para catálogo de negócios locais com cadastro, listagem, avaliações, comentários e informações estruturadas.}
  \resumeItem{Consolida conceitos de CRUD, autenticação, banco de dados e interface responsiva em um fluxo único.}
  \resumeItem{Disponibiliza repositório público e página dedicada no blog para documentação do projeto.}
\resumeItemListEnd

\resumeProjectHeading{\href{https://kaaylooon.github.io/projetos/kaappli/}{\textbf{Kaappli}} | \emph{JavaScript, Electron, HTML/CSS}}{2024}
\resumeItemListStart
  \resumeItem{Aplicativo desktop interativo para prática de Hiragana e Katakana com geração dinâmica de caracteres.}
  \resumeItem{Valida respostas romanizadas, exibe feedback imediato e acompanha progresso de estudo.}
\resumeItemListEnd

\resumeProjectHeading{\href{https://kaaylooon.github.io/projetos/kapraxis/}{\textbf{Kapraxis}} | \emph{C++, Qt}}{2024 -- Presente}
\resumeItemListStart
  \resumeItem{Aplicativo de estudos com flashcards, revisão espaçada, temporizador e estatísticas de desempenho.}
  \resumeItem{Centraliza ferramentas de rotina acadêmica em uma interface desktop única.}
\resumeItemListEnd

\resumeProjectHeading{\href{https://kaaylooon.github.io/projetos/kext/}{\textbf{Kext}} | \emph{C, Terminal UI}}{2024 -- Presente}
\resumeItemListStart
  \resumeItem{Editor de texto inspirado no vi, escrito em C puro com foco em leveza e rapidez no terminal.}
  \resumeItem{Projeto em andamento com modos de inserção e comando, mirando syntax highlighting e múltiplos buffers.}
\resumeItemListEnd

\resumeSubHeadingListEnd

\section{Competências Técnicas}
\begin{itemize}[leftmargin=0.15in, label={}]
\small{\item{
 \textbf{Linguagens}{: Python, TypeScript, JavaScript, SQL} \\
 \textbf{Mobile e Web}{: React Native, Expo, Flask, Django, HTML/CSS, Bootstrap} \\
 \textbf{Backend}{: Node.js} \\
 \textbf{Desktop e Sistemas}{: Qt, Electron, Raylib, Terminal UI (TUI)} \\
 \textbf{Bancos de Dados e Ferramentas}{: PostgreSQL, SQLite, Git, Linux, LaTeX, AsyncStorage} \\
}}
\end{itemize}


\section{Habilidades Extras}
\begin{itemize}[leftmargin=0.15in, label={}]
\small{\item{
 \textbf{Idiomas:} Português nativo; Inglês B2/C1; Francês e Espanhol em início de aprendizado. \\

 \textbf{Base Acadêmica:} Estudo aprofundado de matemática e física, incluindo cálculo, álgebra linear, teoria dos números, relatividade e tópicos de física moderna. \\

 \textbf{Olimpíadas Científicas:} Ouro na OLITEF (2024), O2 (2024) e OBLI (2024); prata na OBA (2025 e 2026); bronze na OBMEP (2025), OBLI (2025), OBG (2024) e OLITEF (2025); menção na OBF (2025). \\
}}
\end{itemize}


\end{document}
